% Generated draft: requires mathematical and source review.
\documentclass[11pt]{article}
\usepackage[margin=25mm]{geometry}
\usepackage{fontspec}
\setmainfont{Latin Modern Roman}
\usepackage{amsmath,amssymb,mathtools,unicode-math}
\setmathfont{Latin Modern Math}
\usepackage{enumitem,needspace,xurl,hyperref}
\setlength{\parindent}{0pt}
\setlength{\parskip}{5pt}
\setlength{\emergencystretch}{3em}
\newcommand{\N}{\mathbb{N}}
\newcommand{\Z}{\mathbb{Z}}
\newcommand{\Q}{\mathbb{Q}}
\newcommand{\R}{\mathbb{R}}
\newcommand{\C}{\mathbb{C}}
\newcommand{\F}{\mathbb{F}}
\newcommand{\A}{\mathbb{A}}
\newcommand{\PP}{\mathbb{P}}
\newcommand{\E}{\mathbb{E}}
\newcommand{\eps}{\varepsilon}
\newcommand{\dd}{\,\mathrm{d}}
\begin{document}
\section*{UnsolvedMath ID: 30006447}
\textbf{Source problem code:} catalogue entry
\subsection{Definitions and mathematical statement}
Let \(f\) denote the source's defining equation for a reduced plane curve with distinguished point \(0\). The base field, the algebraic or analytic category, and whether the objects are germs or global curves are not specified in the supplied record. These choices remain unresolved hypotheses.

For a specified marked curve \(C\) defined by \(f\), the punctual Hilbert scheme \(\operatorname{Hilb}_n(f)_0\) parametrizes zero-dimensional subschemes of \(C\) of length \(n\) supported at \(0\). This describes the terminology in the cleaned statement; the damaged original excerpt instead says ``principal Hilbert schemes,'' and does not provide a reliable additional principal-ideal condition.

Write \(X_{f,n}^0\) for the objects called restricted contact loci in the source. Their precise definition, ambient jet or arc space, restriction, and indexing convention are not supplied. In particular, the superscript \(0\) cannot safely be interpreted from this excerpt alone.

Write \(Q_f(s,T)\) for the source's motivic Igusa series, where \(s\) is a formal parameter and \(T\) is a series variable. The damaged displayed definition visibly involves a sum indexed by \(n\geq 0\), a motivic integral, an order along a discriminant \(\Delta_f\), an orbifold volume-form symbol \(\omega_{\mathrm{orb}}\), and formal powers involving \(\mathbb L^{-s}\) and \(\mathbb L^{1/2}\). Here \(\mathbb L\) conventionally denotes the Lefschetz class. The integration spaces, measure, normalization, discriminant, and coefficient ring cannot be recovered reliably from the supplied text. Thus this draft retains \(Q_f\) as a source-defined symbol rather than asserting a reconstructed integral formula.

\begin{enumerate}
\item \textbf{Source Question 3.} Compute \(Q_f(s,T)\) for every smooth curve \(f\) in the intended setting; the eventual target is to compute it for every reduced plane curve in that setting.
\item \textbf{Source Question 4.} For every reduced plane curve \(f\) in the intended setting, find an ``exponential'' relation between the cohomology of its restricted contact loci \(X_{f,n}^0\) and its punctual Hilbert schemes \(\operatorname{Hilb}_n(f)_0\). The cohomology theory, coefficients, grading, range of indices, generating series, and exponential operation are not specified, so no particular equality is asserted here.
\end{enumerate}

The surrounding excerpt motivates \(Q_f\) by a proposed comparison involving \(P(T)\), an exponential operation, \(Z_f\), and half-powers of \(\mathbb L\); its damaged typography does not determine the exact comparison. It also states that even the smooth-curve computation amounts to computing naive Igusa zeta functions of discriminants of monic degree-\(n\) polynomials for all \(n\). These are supplied contextual claims, not additional verified results or a solution. A fully specified mathematical formulation requires recovery of the missing definitions and conventions from the cited source.
\subsection{Short English statement}
Compute the source-defined motivic Igusa series first for smooth curves and ultimately for all reduced plane curves. Also seek an exponential relation connecting restricted contact-locus cohomology with punctual Hilbert schemes. The supplied excerpt does not determine the series, contact loci, or exponential operation precisely enough to state a unique formula. This is a statement-only draft requiring human mathematical review.
\subsection{Sources}
\begin{itemize}
\item Local UnsolvedMath registry: supplied record with permanent ID \(30006447\), title ``Motivic Igusa Series of Plane Curves,'' exact source problem code \texttt{catalogue entry}. The record attributes the questions to Oberwolfach Reports workshop discussions, records proposed year \(2025\), and names no individual proposer.
\item Supplied source citation: \emph{Singularities (2026)}, Oberwolfach Reports, pp.~2413--2490, DOI \texttt{10.4171/owr/2025/45}: \url{https://doi.org/10.4171/owr/2025/45}. The supplied literature list labels this ``Oberwolfach Report 45/2025.'' The year discrepancy is unresolved.
\item Supplied literature source: \emph{Discriminants and motivic integration}, \url{https://arxiv.org/abs/2503.18449}. No authors or precise supporting passage are supplied.
\item Literature assessment dated exactly \texttt{2026-08-22}: the supplied curation labels the problem open and reports that the smooth-curve case required a family of naive Igusa-zeta computations not then available in the literature, with both the series computation and exponential relation left as questions. This is supplied dated curation, not an independently verified literature assessment. The background explicitly requests independent verification of time-sensitive status and cited claims.
\item The excerpt attributes related work in the \(p\)-adic setting to ``G-Wei-Yin [10]'' and says it was recently used to settle a conjecture by Bhargava--Cremona--Fisher--Gajovi\'{c} on factorization densities of \(p\)-adic polynomials. The bibliographic entry for reference [10] is not supplied; this attribution and historical claim have not been independently verified.
\end{itemize}
\end{document}
